% Do not modify these
\documentclass[fleqn,10pt]{wlscirep}
\usepackage[utf8]{inputenc}
\usepackage[T1]{fontenc}



% -- Insert any custom LaTeX packages here --

% \package{natbib} % <-- Required for the Chicago citation style
% \package{apacite} % <-- Required for the APA citation style
% If you decide to use one of the styles above, remember to change the \bibliographystyle{} at the bottom of the document too!

\usepackage{listings} % <-- Required if you want to display program source code in your paper.
\usepackage{xcolor}
 
\definecolor{codegreen}{rgb}{0,0.6,0}
\definecolor{codegray}{rgb}{0.5,0.5,0.5}
\definecolor{codepurple}{rgb}{0.58,0,0.82}
\definecolor{backcolour}{rgb}{0.95,0.95,0.92}
 
\lstdefinestyle{mystyle}{
    backgroundcolor=\color{backcolour},   
    commentstyle=\color{codegreen},
    keywordstyle=\color{magenta},
    numberstyle=\tiny\color{codegray},
    stringstyle=\color{codepurple},
    basicstyle=\ttfamily\footnotesize,
    breakatwhitespace=false,         
    breaklines=true,                 
    captionpos=b,                    
    keepspaces=true,                 
    numbers=left,                    
    numbersep=5pt,                  
    showspaces=false,                
    showstringspaces=false,
    showtabs=false,                  
    tabsize=2
    }
\usepackage{hyperref}
\hypersetup{
    colorlinks=true,
    linkcolor=blue,
    filecolor=magenta,      
    urlcolor=cyan,
}
\urlstyle{same}

% -- End of custom LaTeX packages --


% Fill in your title
\title{Exploring Bitcoin }

% Do not modify the author tag below, just let it be blank
\author{}

% Fill in assignment abstract
\begin{abstract}
In this assignment we parsed a dataset, in the form of XMLs, through a Hadoop- and Pig apache. We learned how to write MapReduce jobs and Pig scripts, and how effective it is on our dataset based on XML files from Bitcoin stack exchange. 

\end{abstract}


% Do not modify the following two lines
\lstset{style=mystyle}
\begin{document}
% define the infrastructure
\ExplSyntaxOn
\keys_define:nn { coverpage }
 {
  assignment_title .tl_set:N 	= \hk_variable_assignment_title_tl,
  student_name  .tl_set:N 		= \hk_variable_student_name_tl,
  student_number     .tl_set:N 	= \hk_variable_student_number_tl,
  course_name .tl_set:N 			= \hk_variable_course_name_tl,
  course_code .tl_set:N			= \hk_variable_course_code_tl,
  due_date .tl_set:N 			= \hk_variable_due_date_tl,
  master_of .tl_set:N 			= \hk_variable_master_of_tl,
  group_size .tl_set:N 			= \hk_variable_group_size_tl,
 }
\cs_new_protected:Nn \use_coverpage:n
 {
  \keys_set:nn { coverpage } { #1 }
  \coverpage
 }
\cs_set_eq:NN \makecoverpage \use_coverpage:n
\cs_new:Nn \use_variable:n
 {
  \tl_use:c { hk_variable_#1_tl }
 }
\cs_set_eq:NN \use \use_variable:n
\ExplSyntaxOff



% Command for making the cover page

\newcommand{\coverpage}{%
\begin{titlepage}
\begin{center}
\begin{figure}[ht]
\centering
\includegraphics[width=0.5\linewidth]{images/hk-logo.pdf}
\end{figure}
\huge{Course Assignment}\\
\LARGE{Master of \use{master_of}}\\
\Large{Department of Technology}\\

\vskip10pt%
\begin{table}[ht]
\centering
\Large
\setlength{\tabcolsep}{0.5em}
\setlength\extrarowheight{1.2em}
\begin{tabularx}{0.8\linewidth}{lX}
\toprule
 Assignment title 	& \use{assignment_title} \\ \hline
 Course code 		& \use{course_code} \\ \hline
 Course name 		& \use{course_name} \\ \hline
 Due date			& \use{due_date} \\
 \bottomrule
\end{tabularx}%
\end{table}

\vskip10pt
\textbf{\Large{Declaration:}}\\

\if\use{group_size}1
Through the submission of this assignment, I hereby declare\\ that this report is the result of my own work, and that all\\ sources have been properly cited to throughout the text.
\vskip10pt%
\begin{table}[ht]
\centering
\Large
\setlength{\tabcolsep}{0.5em}
\setlength\extrarowheight{1em}
\begin{tabularx}{0.8\linewidth}{lX}
\toprule
 Student name 		& \use{student_name} \\ \hline
 Student ID number 	& \use{student_number} \\ 
 \bottomrule
\end{tabularx}%
\end{table}

\else 

Through the submission of this assignment, we hereby declare\\ that this report is the result of our own work, and that all\\ sources have been properly cited to throughout the text.
\begin{table}[ht]
\centering
\Large
\setlength{\tabcolsep}{0.5em}
\setlength\extrarowheight{1em}
\begin{tabularx}{0.8\linewidth}{lX}
\toprule
 Names of students 		& \use{student_name} \\ \hline
 Student ID numbers 	& \use{student_number} \\ 
 \bottomrule
\end{tabularx}%
\end{table}
\fi




\end{center}
\end{titlepage}
}


% Insert data for the hand-in's cover page
\makecoverpage{
	master_of 		 = \par{Applied Computer Science},  % Use either: Applied Computer Science | Human-Computer Interaction
	assignment_title = \par{ Exploring Bitcoin} ,  % Title of your assignment
	course_code    	 = \par{MA120},  % Course code (ex. MA110)
	course_name      = \par{Big Data},  % Course name (ex. Systems Development)
	due_date		 = \par{18 Oktober},  % Due date
	student_name     = \par{Theodore Midtbø Alstad ; Howie Chen},  % Your name (or names, if group – separate names with ; semicolon)
	student_number   = \par{865317 ; 866354},  % Your student ID number (or numbers, if group – separate ID numbers with ; semicolon)
	group_size		 = 2, % Number of group members (used for the declaration text)
}


% Do not modify the following two lines
\flushbottom
\maketitle


% --INTRODUCTION--
\section*{Introduction}
We choose to work together because both of us have python background, therefor we choose python as main programming language. We explored Bitcoin as given dataset from \url{archive.org/download/stackexchange/bitcoin.stackexchange.com.7z}  through Apaches: Hadoop and Pig. 

% --FUNCTIOS %-- 
\section*{Main functions}
The datasets contains only XML files, which it means values to different attributes have a lot of  ascii characters, punctiations, numbers and HTML tags. We also needed a function to parse the XML files to Hadoop and to create a mapper for the data.  Therefor we created three different main functions thats being reused through the project. : 
\subsection*{CleanBody}

In cleanBody function it formats strings for mapper function. The function makes non-case sensitivity, removes ascii characters , HTML formatting, and treats anything saperated by blanked space or as separated words. Most of the words after separating/splitting will be counted as separate words, and this will effect the results of some tasks. For instance  the name "Jens-Petter" will be counted as two words; "Jens" and "Petter". The input for the function is a string body which will be formatted and it returns a list of formatted string.

\begin{lstlisting}[language=Python, caption=CleanBody function]
def cleanBody(body):
    body = body.lower()
    body = ascii(body)
    body = sub("<.+?>","",body)
    body = body.replace("/", " ")
    body = body.strip()

    for i in ignore_char:
        	body = body.replace(i, "")

    body = body.split(" ")

    return body
\end{lstlisting}
\subsection*{Mapper Core}
mapper\_core is the core of the mapper function; it prints out the relevant data in a format parseable by Hadoop. It functions through three modes, as determined by the parameter mode: "single", "double", and "triple".

\begin{itemize}
  \item Single: Assumes input "words" is a list of words to print. Prints word in words as (word, 1). Ignores empty strings and spaces.
  \item Double: Assumes input "words" is two nested lists to print. Prints word, count in words as (word, count). Ignores empty strings and spaces.
  \item Triple: Assumes input "words" is three nested lists to print. Prints id, score, title in words as (id, score, title). does not ignore empty strings and spaces.
\end{itemize}

\begin{lstlisting}[language=Python, caption=mapper\_core function]
def mapper_core(words, mode="single"):
    if mode == "single":
        for word in words:
            if word not in ["", " "]:
              print("%s %s" %(word,1)) #Emit the word

    elif mode == "double":
        in1, in2 = words
        for word, count in zip(in1,in2):
            if word not in ["", " "]:
                print("%s %s" %(word,count)) #emit the words

    elif mode == "triple":
        in1, in2,in3  = words
        for id, score, title in zip(in1,in2,in3):
            print("%s %s %s" %(id,score, title)) #emit the words


\end{lstlisting}

\subsection*{Xmlparser}
The chosen method to interpret the dataset is to parse to the mapper function. Although parsing an entire XML-file takes up significant memory, this method fits our dataset. It has been separated out as a function so it may be easily replaced by other methods more fit for large files. The input for this function is a XML-file and the output is a parsed XML-file

\begin{lstlisting}[language=Python, caption=xmlparser function]
def xmlparser(infile):
    if not isinstance(infile, str):
        infile = infile.detach()
    mytree = ET.parse(infile)
    myroot = mytree.getroot()
    return myroot
\end{lstlisting}
% --TASK1--
\section*{Task 1 Warmup}
%This part of the task is about get to know how hadoop- and pig apache works and how to parse XML  files throught python code. 
\subsection*{1a) WordCount}
\textbf{Assumption}: Count the words in body of questions PostTypeID="1"  in the \textit{Posts.xml}  file. The result should be how many times a word occur in the body of questions.\\ \\
\textbf{Implementation}The will  \\ \\
\textbf{Notes/Reflection} bye bye  \\ \\

\lstinputlisting[firstline=500,lastline=510,title=1a\_output]{output1a.txt}

\subsection*{1b) Unique words}
\textbf{Assumption}: Write a MapReduce job where the result should be unique words in the titles PostTypeID="1" in the \textit{Posts.xml} File. \\ \\
\textbf{Implementation}  \\ \\
\textbf{Notes/Reflection}
\lstinputlisting[firstline=500,lastline=510,title=1b\_output]{output1b.txt}

\subsection*{1c) MoreThan10}
\textbf{Assumption}: Write a simple python code to check the title length in \textit{Posts.xml}. The result should output how many titles have more than 10 words.  \\ \\
\textbf{Implementation}  \\ \\
\textbf{Notes/Reflection}
\lstinputlisting[firstline=1,lastline=1,title=1c\_output]{output1c.txt}

\subsection*{1d) Stopwords}
\textbf{Assumption}: Write a simple python code based on task 1a to exclude \href{https://raw.githubusercontent.com/naimdjon/stopwords/master/stopwords.txt}{stopwords} from body of questions PostTypeID="1" in the \textit{Posts.xml}. The output should be text file without any stopwords. \\ \\
\textbf{Implementation}  \\ \\
\textbf{Notes/Reflection}
\lstinputlisting[firstline=500,lastline=510,title=1d\_output]{output1d.txt}


\subsection*{1e) Pig top 10}
\textbf{Assumption}: Write a pig script to select top 10 listed words after removing the stopwords from \textit{Posts.xml}. The output should print out top 10 listed words and the corresponding occurrence rate.\\ \\
\textbf{Implementation}  \\ \\
\textbf{Notes/Reflection}
\lstinputlisting[firstline=1,lastline=10,title=1e\_output]{output1e.txt}

\subsection*{1f) Tags}
\textbf{Assumption}: Write a MapReduce job to create a dictionary over unique tags in  \textit{Posts.xml}. The result should print the unique tags.\\ \\
\textbf{Implementation}  \\ \\
\textbf{Notes/Reflection}
\lstinputlisting[firstline=500,lastline=510,title=1f\_output]{output1f.txt}


% --TASK2--
\section*{Task 2 Discover}
%This part of the task is about to looking through several 

\subsection*{2a) Counting }
\textbf{Assumption}: We chose to write a MapReduce job to count the total unique users there are in \textit{Users.xml}. The result will print out how many unique users there are.\\ \\
\textbf{Implementation}  \\ \\
\textbf{Notes/Reflection}
\lstinputlisting[firstline=1,lastline=1,title=2a\_output]{output2a.txt}

\subsection*{2b) Unique users }
\textbf{Assumption}: Write a MapReduce job based on task 2a) to create a mapper  and a reducer functions in \textit{Users.xml}. The result should contain unique users in the dataset.  \\ \\
\textbf{Implementation}  \\ \\
\textbf{Notes/Reflection}
\lstinputlisting[firstline=500,lastline=510,title=2b\_output]{output2b.txt}

\subsection*{2c) Top miners }
\textbf{Assumption}:Write a MapReduce job to find top 10 users based on attribute Reputation="x" in \textit{Users.xml}. The result will print out top 10 users based on their reputation. \\ \\
\textbf{Implementation}  \\ \\
\textbf{Notes/Reflection}
\lstinputlisting[firstline=1,lastline=10,title=2c\_output]{output2c.txt}

\subsection*{2d Top) questions }
\textbf{Assumption}: Write a MapReduce job to find top 10 title questions PostTypeID="1"  based on attribute Score="x" in \textit{Posts.xml}.The result lists top 10 questions using id, question and the score.\\ \\
\textbf{Implementation}  \\ \\
\textbf{Notes/Reflection}


\subsection*{2e) Favorite questions }
\textbf{Assumption}: Write a MapReduce job to find top 10 title questions PostTypeID="1"  based on attribute FavoriteCount="x" in \textit{Posts.xml}.The result lists top 10 questions like id, question and the score.\\ \\ \\ \\
\textbf{Implementation}  \\ \\
\textbf{Notes/Reflection}


\subsection*{2f) Average answers}
\textbf{Assumption}: \\ \\
\textbf{Implementation}  \\ \\
\textbf{Notes/Reflection}


\subsection*{2g) Countries }
\textbf{Assumption}: We chose to write a MapReduce job to discover users by countries in \textit{Users.xml}. The result lists different countries and corresponding users. \\ \\
\textbf{Implementation}  \\ \\
\textbf{Notes/Reflection}


\subsection*{2h) Names}
\textbf{Assumption}We chose to write a MapReduce job to find the most popular names in \textit{Users.xml}. The result lists top 10 common names and how many.  \\ \\
\textbf{Implementation}  \\ \\
\textbf{Notes/Reflection}


\subsection*{2i) Answers }
\textbf{Assumption}: Write a simple python code to find how many titles of questions PostTypeId="x" have at least one answers based on attribute AnswerCount in \textit{Users.xml}. The result prints out how many questions have been answered.  \\ \\
\textbf{Implementation}  \\ \\
\textbf{Notes/Reflection}


% --TASK3--
\section*{Task 3 Numbers}

\subsection*{3a) Bigram }
\textbf{Assumption}: We chose to write a MapReduce job to find the most common pair of adjacent words in \textit{Posts.xml}.For instance, "big data" or "Fast car" are examples of bigram. The result print the most common bigram   \\\\
\textbf{Implementation}  \\ \\
\textbf{Notes/Reflection}


\subsection*{3b) Trigram }
\textbf{Assumption}: This task is based on 3a, to find three words that appear consecutively in \textit{Posts.xml}. The result print the most common trigram.  \\ \\
\textbf{Implementation}  \\ \\
\textbf{Notes/Reflection}


\subsection*{3c) Combiner }
\textbf{Assumption}: Write a MapReduce job and add a combiner before  the data sent to the reducer. The result should be the same as a reducer, but the reducer recive a smaller volume of data.  \\ \\
\textbf{Implementation}  \\ \\
\textbf{Notes/Reflection} Having a combiner in A MapReduce jobs saves us bandwidth and computational strain by decreasing the volume of data sent from the mapper.  What combiner means is to have a semi-reducer after the mapper before sending it to the reducer. 


\subsection*{3d) Useless}
\textbf{Assumption}:  We chose write a MapReduce job  to find how many times the word "useless" in \textit{Posts.xml} occurs in the body of questions PostTypeId="1". The result print how many times the word useless occurs.  \\ \\
\textbf{Implementation}  \\ \\
\textbf{Notes/Reflection}

% --TASK4--
\section*{Task 4 Search engine}
\subsection*{4a) Title index}
\textbf{Assumption}: We chose to write MapReduce job to create index over titles, bodies and answers of questions in \textit{Posts.xml}. We are after having a simple index that lists publications in which a search term/s and occur/s. The result is a list of words and their index appearance in posts \\ \\
\textbf{Implementation}  \\ \\
\textbf{Notes/Reflection}


\section*{Conclusion}
% Do not modify this last lines
\end{document}